\documentclass[11pt]{article}

\usepackage[margin=1in]{geometry}
\usepackage{amsmath,amssymb,amsthm,mathtools}
\usepackage{booktabs}
\usepackage{enumitem}
\usepackage[colorlinks=true,citecolor=blue!50!black]{hyperref}

% ---------- macros ----------
\newcommand{\R}{\mathbb{R}}
\newcommand{\N}{\mathbb{N}}
\newcommand{\Pol}[2]{\Pi_{#1}^{#2}}
\newcommand{\norm}[1]{\lVert #1\rVert}
\newcommand{\smin}{\sigma_{\min}}
\newcommand{\rank}{\operatorname{rank}}
\newcommand{\supp}{\operatorname{supp}}
\newcommand{\diag}{\operatorname{diag}}
\newcommand{\pinv}{\dagger}
\DeclareMathOperator*{\argmin}{arg\,min}

\theoremstyle{plain}
\newtheorem{theorem}{Theorem}
\newtheorem{proposition}[theorem]{Proposition}
\theoremstyle{definition}
\newtheorem{hyp}{Hypothesis}
\theoremstyle{remark}
\newtheorem{remark}{Remark}

\title{Weight growth for least-squares RBF-FD stencils with fixed auxiliary centers:\\ an $O(h^{-k})$ bound}
\author{}
\date{}

\begin{document}
\maketitle

\begin{abstract}
We consider RBF-FD weights obtained by least squares against kernel translates centered at auxiliary points that are fixed as the stencil spacing $h\to0$.
By comparing the least-squares weights with the polynomial-exact finite-difference weights, we show $\norm{w_h}_1\le C h^{-k}$ for an operator of order $k$.
The argument covers square stencils ($N=\dim\Pi_p^d$) and, under one additional rank condition on the nodes, arbitrary larger stencils.
When that rank condition fails, the exponent degrades in a precise way.
\end{abstract}

\section{Setting and notation}

Let $z\in\R^d$ be the target point and
$L=\sum_{|\alpha|\le k}a_\alpha\partial^\alpha$ a differential operator of order $k$ whose coefficients are frozen at $z$.
For a multi-index $\alpha\in\N_0^d$ write $|\alpha|$, $\alpha!$, $\xi^\alpha$ as usual, and set
\[
n_t:=\dim\Pol{t}{d}=\binom{t+d}{d}.
\]

\paragraph{Stencil and centers.}
The stencil has $N$ nodes $x_j=z+h\xi_j$, $j=1,\dots,N$, with reference nodes $\xi_j$ independent of $h$.
The auxiliary centers $c_1,\dots,c_m$ are \emph{fixed independently of $h$}.
For a kernel $K$ (derivatives taken in the second argument) define
\[
A_h\in\R^{m\times N},\quad (A_h)_{\ell j}=K(c_\ell,x_j),
\qquad
b\in\R^m,\quad b_\ell=(L_yK)(c_\ell,z),
\]
and the least-squares weights
\[
w_h=\argmin_{w\in\R^N}\norm{A_hw-b}_2 .
\]
In functional terms, $\sum_j w_j\delta_{x_j}\approx\delta_z\circ L$ on the test functions $K(c_\ell,\cdot)$.

\paragraph{Taylor objects.}
For each degree $t\ge0$, with columns indexed by $|\alpha|\le t$, set
\begin{align*}
(V_t)_{j\alpha}&=\xi_j^\alpha\in\R^{N\times n_t}, &
D_{h,t}&=\diag\big(h^{|\alpha|}\big)\in\R^{n_t\times n_t},\\
(J_t)_{\ell\alpha}&=\tfrac{1}{\alpha!}\,\partial_y^\alpha K(c_\ell,z)\in\R^{m\times n_t}, &
R_{h,t}&=A_h-J_tD_{h,t}V_t^{\top}.
\end{align*}
Let $p$ be the polynomial-reproduction degree and define the operator vector $\eta_L\in\R^{n_p}$ by
\[
(\eta_L)_\alpha=\begin{cases}\alpha!\,a_\alpha, & |\alpha|\le k,\\ 0,& k<|\alpha|\le p.\end{cases}
\]
(For $\alpha=e_i+e_j$, $i\ne j$, the coefficient $a_\alpha$ collects \emph{both} $(i,j)$ and $(j,i)$ terms; e.g.\ for $Lu=-\nabla\!\cdot(A\nabla u)$ one has $a_{e_i+e_j}=-(A_{ij}+A_{ji})$, $a_{2e_i}=-A_{ii}$, and $a_{e_j}=-\sum_i\partial_iA_{ij}$.)

\paragraph{Two degrees.}
\[
p:=\max\{t:\ n_t\le N\},\qquad s:=\min\{t:\ \rank V_t=N\}.
\]
We assume $p\ge k$ (equivalently $N\ge n_k$).
Always $s\ge p$, with $s=p$ if and only if $N=n_p$; for generic nodes with $N>n_p$ one has $s=p+1$.
For distinct nodes $s\le N-1$, so $s$ is always finite.

\section{Hypotheses and main result}

\begin{hyp}[smoothness]\label{h:smooth}
$y\mapsto K(c_\ell,y)$ is $C^{s+1}$ on a ball around $z$ containing all stencils with $h\le h_0$, where $h_0\le 1$.
\end{hyp}
\begin{hyp}[polynomial unisolvence]\label{h:poly}
$\rank V_p=n_p$.  (For $N=n_p$ this says $V_p$ is nonsingular.)
\end{hyp}
\begin{hyp}[centers]\label{h:centers}
$J_s$ has full column rank, i.e.\ $\gamma:=\smin(J_s)>0$. In particular $m\ge n_s$.
\end{hyp}

Define the minimum-norm polynomial-exact weights
\[
w_h^*:=(V_p^{\top})^{\pinv}D_{h,p}^{-1}\eta_L .
\]
They satisfy $D_{h,p}V_p^\top w_h^*=\eta_L$, which is equivalent to
$\sum_j w_{h,j}^*\,q(x_j)=(Lq)(z)$ for all $q\in\Pol{p}{d}$.

\begin{theorem}\label{thm:main}
Under Hypotheses~\ref{h:smooth}--\ref{h:centers} there are $h_1\in(0,1]$ and constants $C,C'$ independent of $h$ such that for $0<h\le h_1$ the matrix $A_h$ is injective (so $w_h$ is unique) and
\begin{equation}\label{eq:main}
\norm{w_h-w_h^*}_2\le C\,h^{\,p+1-s}\,\norm{w_h^*}_2,
\qquad
\norm{w_h^*}_2\le \frac{\norm{\eta_L}_2}{\smin(V_p)}\,h^{-k}.
\end{equation}
Consequently
\[
\norm{w_h}_1\le\sqrt N\,\norm{w_h}_2\le C'\,h^{-k-(s-p-1)_+}.
\]
In particular $\norm{w_h}_1\le C'h^{-k}$ whenever $s\le p+1$.
\end{theorem}

For $N=n_p$ ($s=p$) the least-squares weights are a relative $O(h)$ perturbation of the classical finite-difference weights.
For generic $N>n_p$ ($s=p+1$) they have the same order but are not asymptotically equal to $w_h^*$.

\section{Proof}

\paragraph{Step 1 (Taylor).}
For $t\in\{p,s\}$, Taylor's theorem with Lagrange remainder, applied entrywise ($m,N$ fixed), gives
\[
A_h=J_tD_{h,t}V_t^\top+R_{h,t},\qquad \norm{R_{h,t}}_2\le C_R\,h^{t+1}\quad(0<h\le h_0).
\]

\paragraph{Step 2 (stability).}
For $h\le 1$ we have $D_{h,s}\succeq h^sI$. By Hypothesis~\ref{h:centers} and the definition of $s$, with $v:=\smin(V_s)=\min_{\norm w=1}\norm{V_s^\top w}>0$,
\[
\norm{A_hw}\ge\norm{J_sD_{h,s}V_s^\top w}-\norm{R_{h,s}w}
\ge\big(\gamma\,v\,h^s-C_Rh^{s+1}\big)\norm w
\ge\tfrac12\gamma v\,h^s\norm w
\]
for $h\le h_1:=\min\{h_0,\gamma v/(2C_R)\}$. Hence $A_h$ is injective and
\[
\norm{A_h^\pinv}_2\le \frac{2}{\gamma v}\,h^{-s}.
\]

\paragraph{Step 3 (consistency).}
Since $L$ has order $k\le p$, the identity $b=J_p\eta_L$ holds \emph{exactly}.
Using $D_{h,p}V_p^\top w_h^*=\eta_L$,
\[
A_hw_h^*-b=J_p\big(D_{h,p}V_p^\top w_h^*-\eta_L\big)+R_{h,p}w_h^*=R_{h,p}\,w_h^* .
\]
Moreover, $\eta_L$ is supported on $|\alpha|\le k$ and $h\le1$, so $\norm{D_{h,p}^{-1}\eta_L}\le h^{-k}\norm{\eta_L}$, and therefore
$\norm{w_h^*}\le \smin(V_p)^{-1}\norm{\eta_L}\,h^{-k}$.
This is the only place where $k$ (rather than $p$) enters; it is why the exponent is $-k$.

\paragraph{Step 4 (conclusion).}
Injectivity gives $w_h=A_h^\pinv b$ and $A_h^\pinv A_h=I$. Hence
\[
w_h-w_h^*=A_h^\pinv\big(b-A_hw_h^*\big)=-A_h^\pinv R_{h,p}\,w_h^*,
\]
and by Steps 1--3,
\[
\norm{w_h-w_h^*}\le\frac{2}{\gamma v}h^{-s}\cdot C_Rh^{p+1}\norm{w_h^*}.
\]
This is \eqref{eq:main}. If $s\le p+1$ then $h^{p+1-s}\le1$ and $\norm{w_h}\le(1+C)\norm{w_h^*}$; if $s>p+1$ the factor $h^{p+1-s}$ gives the stated exponent. Finally $\norm w_1\le\sqrt N\norm w_2$. \hfill$\qquad\blacksquare$

\section{Remarks}

\begin{remark}[relation to the $\eta$-formulation]
For $N=n_p$ put $\eta=D_hV^\top w$ and $B_h=A_hV^{-\top}D_h^{-1}=J+R_hV^{-\top}D_h^{-1}$. Then $A_h=B_hD_hV^\top$, $\eta_h=B_h^\pinv b$, and $\eta_h-\eta_L=-B_h^\pinv R_hw_h^*$.
The statement ``$B_h=J+O(h)$'' is the stability estimate of Step~2 (via Weyl's inequality, $\smin(B_h)\ge\gamma/2$).
The statement ``column $\alpha$ of $B_h-J$ is $O(h^{p+1-|\alpha|})$'' is the same as $R_hw_h^*=O(h^{p+1-k})$ of Step~3.
The operator-norm bound $B_h-J=O(h)$ alone yields only $\norm{w_h}\lesssim h^{-p}$: the structure $b=J\eta_L$ is what produces $h^{-k}$.
\end{remark}

\begin{proposition}[sharpness]\label{prop:sharp}
Assume $s\le p+1$ and $a_\beta\ne0$ for some $|\beta|=k$. Then $\norm{w_h}_1\ge c\,h^{-k}$ for small $h$.
\end{proposition}
\begin{proof}
Write $A_hw_h-b=J_p(D_{h,p}V_p^\top w_h-\eta_L)+R_{h,p}w_h$. By least-squares optimality and Step~3, $\norm{A_hw_h-b}\le\norm{A_hw_h^*-b}\le C h^{p+1-k}$; also $\norm{R_{h,p}w_h}\le C_Rh^{p+1}\norm{w_h}\le Ch^{p+1-k}$ by Theorem~\ref{thm:main}.
Since $J_p$ consists of columns of $J_s$, $\smin(J_p)\ge\gamma$, so $\norm{D_{h,p}V_p^\top w_h-\eta_L}\le C\gamma^{-1}h^{p+1-k}$.
For $|\beta|=k$ the $\beta$-th component gives $\big|h^k\sum_j w_{h,j}\xi_j^\beta-\beta!a_\beta\big|\le C h^{p+1-k}$, hence
$h^k\norm{w_h}_1\max_j|\xi_j^\beta|\ge\beta!|a_\beta|-Ch^{p+1-k}$, and $p+1-k\ge1$.
\end{proof}

\begin{remark}[scope]
\begin{enumerate}[leftmargin=*,label=(\roman*)]
\item \emph{Fixed centers.} If the centers move with $h$ (e.g.\ they are the stencil nodes), $J_s$ becomes $h$-dependent and the proof does not carry over; that is the genuine flat-limit regime (cf.\ Schaback 2008, Larsson--Fornberg 2005).
\item \emph{Dependence of constants.} $C$ depends on $\norm{\eta_L}$, $\smin(V_p)$, $\smin(V_s)$, $\gamma$ and on bounds for the derivatives of $K(c_\ell,\cdot)$ of order $s+1$. It is not uniform over center configurations for which $J_s$ degenerates, nor over node configurations for which $V_s$ does.
\item \emph{Number of centers.} $m\ge n_s$ is what this short proof needs. In the experiments below, generic centers with $m=N$ (and $m=N+1$) also gave bounded $h^k\norm{w_h}$; the weaker sufficient condition (injectivity of a graded limit matrix, requiring only $m\ge N$) is not proved here.
\item \emph{Stencil sizes.} In $d=2$ the admissible square sizes are $n_t=21,28,36,45,55,\dots$; in $d=3$ they are $20,35,56,\dots$. Stencils of size $25$ or $50$ therefore fall in the rectangular case: $p=5,8$ in $d=2$ ($p=3,4$ in $d=3$), and one must verify $\rank V_{p+1}=N$ and (for this proof) $m\ge n_{p+1}$.
\item \emph{Structured stencils.} The condition $\rank V_{p+1}=N$ can fail for Cartesian blocks and for nodes on few lines or circles. For a $5\times5$ tensor grid (zeros of a quintic in each variable), the quotient dimensions are $22,24,25$ in degrees $6,7,8$, so $s=8$.
\end{enumerate}
\end{remark}

\section{Numerical check}

Gaussian kernel $K(c,y)=e^{-|c-y|^2}$, $d=2$, $L=\Delta$ ($k=2$), $30$--$40$ random centers fixed outside the unit disk, $z=0$, $200$-digit arithmetic. Table~\ref{tab:num} reports $h^2\norm{w_h}_2$ as $h$ decreases from $10^{-2}$ to $10^{-6}$.

\begin{table}[h]
\centering
\begin{tabular}{@{}llll@{}}
\toprule
Stencil & $s$ & $h^2\norm{w_h}_2$ & Observation\\
\midrule
21 nodes, triangular lattice ($N=n_5$) & $5$ & $507.5\to504.23$ & $\to h^2\norm{w_h^*}$; $\norm{w_h-w_h^*}/\norm{w_h^*}\approx0.7h$\\
25 generic nodes & $6$ & $1036\to1004.4$ & bounded; relative gap $\to2.9$ (not $O(h)$)\\
$5\times5$ grid ($\rank V_6=22$) & $8$ & $9.5\!\cdot\!10^{3}\to9.5\!\cdot\!10^{11}$ & $\norm{w_h}\sim h^{-4}=h^{-k-(s-p-1)}$\\
\bottomrule
\end{tabular}
\caption{Scaling of least-squares weights. The third row confirms that the rank condition $s\le p+1$ is necessary for the exponent $-k$.}
\label{tab:num}
\end{table}

The result is asymptotic. For a \emph{random} 21-point stencil in the unit disk, $\norm{w_h^*}\approx1.4\cdot10^{6}h^{-2}$, and the least-squares weights had not entered the asymptotic regime at $h=10^{-2}$: well-conditioned stencils are needed to see the $h^{-k}$ slope at practical $h$.

\section*{Background references}
The following are background for the ingredients (derivative reproducing kernels, coalescing-node limits, polynomial-exact differentiation formulas). None proves the bound above; bibliographic details should be verified before citing.
\begin{itemize}[leftmargin=*]
\item D.-X.\ Zhou, Derivative reproducing properties for kernel methods in learning theory, \emph{J.\ Comput.\ Appl.\ Math.} 220 (2008).
\item R.\ Schaback, Limit problems for interpolation by analytic radial basis functions, \emph{J.\ Comput.\ Appl.\ Math.} 212 (2008).
\item O.\ Davydov and R.\ Schaback, Error bounds for kernel-based numerical differentiation, \emph{Numer.\ Math.} 132 (2016).
\item O.\ Davydov and R.\ Schaback, Minimal numerical differentiation formulas, \emph{Numer.\ Math.} 140 (2018).
\end{itemize}

\end{document}
